\documentclass[leqno]{article}
% Shared mark scheme preamble for resources.danieljsmith.org: the pack
% preamble, unchanged, then the scheme's own layer. A scheme's .tex is
% rendered from its YAML (scripts/README.md, Mark schemes), never edited.
\input{\CurrentFilePath/pack_preamble.tex}

% The page: equal side margins, so the table runs the full width of the
% header rule. The footers' right edge follows the text's.
\geometry{left=0.8in, right=0.8in}
\makeatletter
\renewcommand{\djs@footR}[1]{%
  \put(\dimexpr1in+\hoffset+\oddsidemargin+\textwidth\relax,
    -\dimexpr\paperheight-0.32in\relax){\makebox[0pt][r]{#1}}}
\makeatother

\usepackage[export]{adjustbox}
\usepackage{tabularray}
\definecolor{msGrid}{RGB}{176, 172, 164}
\definecolor{msTint}{RGB}{245, 243, 239}

% The header on every page, the front page included, and the booklet label
% under the contents-page title. \djsQuestionPageStyle stays empty.
\newcommand{\djsSchemeHeader}[1]{%
  \djsPackHeader[markscheme]{\textit{Mark scheme} - #1}%
  \renewcommand{\djsPackTitle}{#1}%
  \renewcommand{\djsBookletLabel}{Mark scheme}%
}

% The front page: the contents heading, spaced as tocloft spaces the packs'
% "Questions", then the contents table. A question's name and page number
% both link to its first page.
\newcommand{\msContentsHeading}{%
  \vskip\cftbeforetoctitleskip
  {\noindent\cfttoctitlefont\contentsname\par}%
  \vskip\cftaftertoctitleskip}
\NewTblrEnviron{mscontents}
\SetTblrInner[mscontents]{
  colspec = {Q[wd=26mm, l] X[l] Q[wd=16mm, c] Q[wd=14mm, r]},
  width = \linewidth,
  rowsep = 4pt, colsep = 6pt,
  hlines = {0.4pt, msGrid}, vlines = {0.4pt, msGrid},
  hline{2} = {2.2pt, scarlet},
  row{1} = {font=\small\bfseries},
  row{Z} = {font=\bfseries},
  column{2} = {font=\small},
}
\newcommand{\msContentsQuestion}[1]{\hyperref[q:#1]{Question #1}}
\newcommand{\msContentsPage}[1]{\hyperref[q:#1]{\pageref*{q:#1}}}
\newcommand{\msSep}{\hspace{0.9em}}

% The conventions pages' layout: headings in the documents' own bold, tables
% ruled and padded as the scheme's, roomy lists.
\newcommand{\msHead}[1]{%
  \par\vspace{18pt}%
  {\noindent\bfseries #1\par}%
  \vspace{6pt}}
\NewTblrEnviron{mstblr}
\SetTblrInner[mstblr]{rowsep = 6pt, colsep = 6pt,
  hlines = {0.4pt, msGrid}, vlines = {0.4pt, msGrid}}
\newenvironment{msList}
  {\begin{itemize}[leftmargin=1.2em, itemsep=5pt, topsep=3pt, parsep=0pt]}
  {\end{itemize}}

% A question's first page: the link target and the PDF bookmark, as
% \questionitem gives the packs.
\newcommand{\msQnum}{}
\newcommand{\msQuestion}[1]{%
  \clearpage
  \renewcommand{\msQnum}{#1}%
  \phantomsection\label{q:#1}%
  \addcontentsline{toc}{subsection}{\protect\djsTocEntry{#1}}}

% The scheme's table: Question 16mm, Scheme the rest, Marks 16mm, Notes
% 44mm; every cell ruled in msGrid; the scarlet rule under the head row. A
% long table has no caption or "continued" foot; a continuation page opens
% with a muted line above the repeated head row. The rules across a part's
% rows, the alternative's tint and the total row are the rendered table's.
\NewTblrTheme{ms}{%
  \DefTblrTemplate{firsthead}{default}{}%
  \DefTblrTemplate{middlehead,lasthead}{default}{%
    {\small\itshape\color{djsMuted}Question \msQnum\ continued\par}%
    \vspace{6pt}}%
  \DefTblrTemplate{firstfoot,middlefoot,lastfoot}{default}{}%
}
\NewTblrEnviron{msscheme}
\SetTblrOuter[msscheme]{long, theme=ms, entry=none, label=none}
\SetTblrInner[msscheme]{
  colspec = {Q[wd=16mm, l] X[l] Q[wd=16mm, c] Q[wd=44mm, l]},
  width = \linewidth,
  rowsep = 6pt, colsep = 6pt,
  rowhead = 1,
  vlines = {0.4pt, msGrid},
  hline{1} = {0.4pt, msGrid}, hline{2} = {2.2pt, scarlet},
  hline{Z} = {0.4pt, msGrid},
  column{1} = {font=\bfseries}, column{3} = {font=\bfseries},
  column{4} = {font=\small},
  row{1} = {rowsep=4pt, font=\small\bfseries},
}

% Two marks on a row: each note starts with its code.
\newcommand{\msNote}[2]{\textbf{#1}\enspace #2}
% A full-row figure: the worked solution's own picture, scaled uniformly and
% only down, to the table less the Question column. The picture is drawn
% with the line width it has in the solutions document, at its list depth
% there (the optional argument, 1 by default: a part's working), so one
% sized by \linewidth keeps its size; tabularray sets the cell's \linewidth
% to the cell's measured width, which would shrink it at every pass. The
% widths are measured once, in a solution box at the packs' text width
% (pack_preamble.tex's margins, 0.8in and 1in).
\newlength{\msSpanW}
\makeatletter
\newcommand{\ms@keepLW}[1]{%
  \expandafter\xdef\csname ms@solnLW@#1\endcsname{\the\linewidth}}
\AtBeginDocument{%
  \setlength{\msSpanW}{\dimexpr\textwidth-16mm-28pt\relax}%
  \setbox\z@\vbox{%
    \begin{minipage}{\dimexpr\paperwidth-1.8in\relax}
      \begin{djsSolution}
        \ms@keepLW{0}%
        \begin{enumerate}\item \ms@keepLW{1}%
        \begin{enumerate}\item \ms@keepLW{2}%
        \begin{enumerate}\item \ms@keepLW{3}%
        \end{enumerate}\end{enumerate}\end{enumerate}
      \end{djsSolution}
    \end{minipage}}}
\newcommand{\msFigure}[2][1]{%
  \adjustbox{max width=\msSpanW, max totalheight=0.6\textheight, valign=t}{%
    \hsize=\csname ms@solnLW@#1\endcsname
    \linewidth=\hsize
    #2}}
\makeatother
% A figure mark's feature, in the Scheme column.
\newcommand{\msFeature}[1]{\textit{#1}}
% An alternative's title row.
\newcommand{\msAlternative}[1]{\textit{Alternative:} #1}
% Drafts only: a queried mark's code, the query's why under its note, and a
% flag in its part's place.
\newcommand{\msQueried}[2]{#1~\(\to\) #2?}
\newcommand{\msQuery}[1]{{\itshape\color{djsMuted}Query: #1}}
\newcommand{\msFlag}[2]{\textbf{#1:}\enspace #2}


\djsSchemeHeader{Linear Transformations}
\djsSitePath{Core Pure}{Vectors, Matrices \& Linear Transformations}

\djsLastUpdated{26 October 2026}
\begin{document}
\djsSetQuestionCount{11}
\djsAddMarks{1}{6}
\djsAddMarks{2}{10}
\djsAddMarks{3}{9}
\djsAddMarks{4}{10}
\djsAddMarks{5}{9}
\djsAddMarks{6}{8}
\djsAddMarks{7}{2}
\djsAddMarks{8}{5}
\djsAddMarks{9}{11}
\djsAddMarks{10}{4}
\djsAddMarks{11}{12}
\vspace*{6pt}
\djsContentsTitle
\vspace{30pt}
\msContentsHeading
\begin{mscontents}{}
Question & Marks by part & Marks & Page \\
\msContentsQuestion{1} & (a)\,1\msSep (b)\,1\msSep (c)\,1\msSep (d)\,1\msSep (e)\,2 & 6 & \msContentsPage{1} \\
\msContentsQuestion{2} & (a)\,2\msSep (b)\,2\msSep (c)\,1\msSep (d)\,2\msSep (e)\,3 & 10 & \msContentsPage{2} \\
\msContentsQuestion{3} & (i)(a)\,1\msSep (i)(b)\,1\msSep (i)(c)\,1\msSep (ii)\,6 & 9 & \msContentsPage{3} \\
\msContentsQuestion{4} & (a)\,6\msSep (b)\,3\msSep (c)\,1 & 10 & \msContentsPage{4} \\
\msContentsQuestion{5} & (i)\,1\msSep (ii)\,1\msSep (iii)\,2\msSep (iv)\,2\msSep (v)\,1\msSep (vi)\,2 & 9 & \msContentsPage{5} \\
\msContentsQuestion{6} & (a)\,2\msSep (b)\,3\msSep (c)\,3 & 8 & \msContentsPage{6} \\
\msContentsQuestion{7} & & 2 & \msContentsPage{7} \\
\msContentsQuestion{8} & (a)\,1\msSep (b)\,2\msSep (c)\,1\msSep (d)\,1 & 5 & \msContentsPage{8} \\
\msContentsQuestion{9} & (a)\,1\msSep (b)\,2\msSep (c)\,1\msSep (d)\,2\msSep (e)\,3\msSep (f)\,2 & 11 & \msContentsPage{9} \\
\msContentsQuestion{10} & (a)\,3\msSep (b)\,1 & 4 & \msContentsPage{10} \\
\msContentsQuestion{11} & (a)\,2\msSep (b)(i)\,1\msSep (b)(ii)\,2\msSep (c)\,6\msSep (d)\,1 & 12 & \msContentsPage{11} \\
\SetCell[c=2]{r} Total & & 86 & \\
\end{mscontents}
\newpage
% The conventions pages, shared by every pack's mark scheme: one source of
% truth, \input after the front page, with no strand or pack logic. The keys
% on the first side, the rules on the second. The scheme layer defines
% \msHead, msList and mstblr.
\vspace*{6pt}%
{\noindent\cfttoctitlefont Reading the mark scheme\par}%
\vspace{10pt}%
These rules apply throughout, unless a note says otherwise.

\msHead{Marks}
Every mark is worth one, so \textbf{M1 A1} on a row is two marks. Only marks
coded ft follow through.\par
\vspace{10pt}%
\begin{mstblr}{
  colspec = {Q[wd=15mm, l] X[l]},
  width = \linewidth,
  hline{2} = {2.2pt, scarlet},
  column{1} = {font=\bfseries},
  row{1} = {font=\small\bfseries, rowsep=4pt},
}
Code & Meaning \\
M1 & Method: a correct method, applied to this question. A slip doesn't cost it if the method is clear. A formula must be used, not just quoted. \\
dM1 & Dependent method: as M1, but only if the method mark before it was earned, unless its note names another. \\
A1 & Accuracy: a correct answer or step. It needs the method mark it follows. \\
A1ft & As A1, following through from an earlier answer that the note names. \\
A1~AG & The last mark of a part whose answer is printed in the question: complete working, with no errors. \\
B1 & A correct value, fact or statement. It needs no method mark. \\
B1ft & As B1, following through from an earlier answer that the note names. \\
E1 & A written explanation, interpretation or conclusion. Concluding a proof by induction needs every earlier mark. \\
\end{mstblr}

\msHead{Abbreviations}
ft and AG qualify a code in the Marks column. The others appear in the
notes.\par
\vspace{10pt}%
\begin{mstblr}{
  colspec = {Q[wd=17mm, l] X[l] Q[wd=15mm, l] X[l]},
  width = \linewidth,
  column{1} = {font=\bfseries}, column{3} = {font=\bfseries},
}
ft & follow through & AG & answer given in the question \\
cao & correct answer only, with no follow-through & cso & correct solution only: no errors in the part \\
oe & or an equivalent form & awrt & any answer that rounds to this value \\
soi & seen, or implied by later working & isw & ignore working after a correct answer \\
SC & special case: credit outside the main method & their & the student's own earlier value \\
condone & a slip that doesn't cost the mark & [\;\,] & working the mark doesn't need \\
s.f. & significant figures & d.p. & decimal places \\
\end{mstblr}

\newpage
\msHead{Reading the scheme}
\begin{msList}
\item Each row is a line of working, with the marks it earns beside it.
  Where a row carries two marks, each note starts with its code.
\item \textbf{[\textit{n}]} ends each part with its total.
\item A shaded \textit{Alternative} is a second method for the same total.
  Mark an answer by the method it follows.
\item Any other valid method earns the equivalent marks, unless the question
  names a method or says \textit{hence}.
\item A sketch earns one B1 for each feature named below its figure.
\end{msList}

\msHead{Accuracy and answers}
\begin{msList}
\item Unless a question says otherwise, give answers to 3 significant
  figures, and angles in degrees to 1 decimal place. Angles in radians take
  3 significant figures.
\item A final answer with more figures is accepted if it rounds to the
  scheme's value; one with fewer is not. An answer spoiled by rounding too
  early loses its accuracy mark.
\item Where a question asks for an exact answer, a decimal loses the final
  accuracy mark. Otherwise either is accepted.
\item Wrong or missing units cost nothing unless the question asks for them.
\item A correct answer with no working earns full marks, unless the question
  needs the working: show that, prove, hence, a method it names, or an exact
  answer. A wrong answer with no working earns nothing. Calculator answers
  are treated the same way.
\item A correct answer, once reached, isn't spoiled by later working (isw).
\end{msList}

\msHead{Value of $\boldsymbol{g}$}
\begin{msList}
\item Take $g = 9.8\,\mathrm{m\,s^{-2}}$ unless a question states a value,
  which is then the only one accepted. Answers that use $g = 9.8$ may be
  given to 2 or 3 significant figures.
\item Where a question doesn't state $g$, an answer worked with $g = 10$
  throughout earns the same marks, keeping any value printed in the
  question. The worked solutions list the final answers with $g = 10$.
\end{msList}


\msQuestion{1}
\begin{msscheme}{
  hline{3,5,7,9,11-12} = {2-Z}{0.4pt, msGrid},
  hline{4,6,8,10,13} = {0.4pt, msGrid},
  row{13} = {font=\bfseries},
}
Question & Scheme & Marks & Notes \\
1(a) & {\(\mathbf{P} = \begin{pmatrix}1&0\\0&-1\end{pmatrix}\)} & B1 & \\
\nopagebreak & & [1] & \\
1(b) & {\(\mathbf{Q} = \begin{pmatrix}0&1\\1&0\end{pmatrix}\)} & B1 & \\
\nopagebreak & & [1] & \\
1(c) & {\(\mathbf{R} = \mathbf{Q}\mathbf{P}\)} & B1 & {B0 for \(\mathbf{P}\mathbf{Q}\).} \\
\nopagebreak & & [1] & \\
1(d) & {\(\mathbf{R} = \begin{pmatrix}0&-1\\1&0\end{pmatrix}\)} & B1ft & {ft their \(\mathbf{P}\) and \(\mathbf{Q}\), multiplied as \(\mathbf{Q}\mathbf{P}\).} \\
\nopagebreak & & [1] & \\
1(e) & {Rotation about the origin} & B1 & {Accept about \(O\) or \((0,0)\).} \\
& {through \(90^\circ\) anticlockwise} & B1 & {Accept \(\frac{\pi}{2}\) anticlockwise or \(270^\circ\) clockwise. Direction required.} \\
\nopagebreak & & [2] & \\
\nopagebreak \SetCell[c=2]{r} Total for Question 1 & & 6 & \\
\end{msscheme}

\msQuestion{2}
\begin{msscheme}{
  hline{3-4,6-7,9,11-12,14-16} = {2-Z}{0.4pt, msGrid},
  hline{5,8,10,13,17} = {0.4pt, msGrid},
  row{17} = {font=\bfseries},
}
Question & Scheme & Marks & Notes \\
2(a) & {\(\begin{pmatrix}0&-1\\1&0\end{pmatrix}\begin{pmatrix}-2&1&1\\1&1&5\end{pmatrix} = \ldots\)} & M1 & {Applies \(\mathbf{P}\) to the vertices, by multiplying or by the rule \((x,y) \mapsto (-y,x)\). One correct image is enough.} \\
& {\(A'(-1,-2)\), \(B'(-1,1)\), \(C'(-5,1)\)} & A1 & {All three, as coordinates or column vectors.} \\
\nopagebreak & & [2] & \\
2(b) & {Rotation about the origin} & B1 & {Needs rotation and the centre \(O\) or \((0,0)\).} \\
& {through \(90^\circ\) anticlockwise} & B1 & {Accept \(\frac{\pi}{2}\) anticlockwise or \(270^\circ\) clockwise. Condone \(+90^\circ\) with no direction.} \\
\nopagebreak & & [2] & \\
2(c) & {\(\mathbf{QR} = \begin{pmatrix}1&4\\3&0\end{pmatrix}\)} & B1 & \\
\nopagebreak & & [1] & \\
2(d) & {\(\det(\mathbf{QR}) = 1\times 0 - 4\times 3\)} & M1 & {Uses \(ad - bc\) on their \(2\times2\) \(\mathbf{QR}\), or multiplies \(\det\mathbf{Q}\) by \(\det\mathbf{R}\). Condone one sign slip.} \\
& {\(\phantom{\det(\mathbf{QR})} = -12\)} & A1 & \\
\nopagebreak & & [2] & \\
2(e) & {\(\text{Area of } T = \tfrac12\times 3\times 4 = 6\)} & B1 & \\
& {\(\text{Area of } T'' = |\text{"}{-12}\text{"}| \times 6\)} & M1 & {Multiplies their area of \(T\) by the modulus of their determinant from (d). Not for an area found from the vertices of \(T''\).} \\
& {\(\phantom{\text{Area of } T''} = 72\)} & A1ft & {ft their area of \(T\) and their determinant.} \\
\nopagebreak & & [3] & \\
\nopagebreak \SetCell[c=2]{r} Total for Question 2 & & 10 & \\
\end{msscheme}

\msQuestion{3}
\begin{msscheme}{
  hline{3,5,7,9-13} = {2-Z}{0.4pt, msGrid},
  hline{4,6,8,14} = {0.4pt, msGrid},
  row{14} = {font=\bfseries},
}
Question & Scheme & Marks & Notes \\
3(i)(a) & {\(\begin{pmatrix}0 & 1\\1 & 0\end{pmatrix}\)} & B1 & \\
\nopagebreak & & [1] & \\
3(i)(b) & {\(\begin{pmatrix}\frac12 & \frac{\sqrt3}{2}\\-\frac{\sqrt3}{2} & \frac12\end{pmatrix}\)} & B1 & {Exact values, or awrt 0.866 for \(\frac{\sqrt3}{2}\).} \\
\nopagebreak & & [1] & \\
3(i)(c) & {\(\begin{pmatrix}0 & 1\\1 & 0\end{pmatrix}\begin{pmatrix}\frac12 & \frac{\sqrt3}{2}\\-\frac{\sqrt3}{2} & \frac12\end{pmatrix} = \begin{pmatrix}-\frac{\sqrt3}{2} & \frac12\\\frac12 & \frac{\sqrt3}{2}\end{pmatrix}\)} & B1ft & {ft their (a) times their (b), in that order. Not for the product in the reverse order.} \\
\nopagebreak & & [1] & \\
3(ii) & {\(\det\begin{pmatrix}2 & 5\\1 & -1\end{pmatrix} = -2 - 5 = -7\)} & B1 & \\
& {Area of \(T = \dfrac{84}{\text{"}7\text{"}} = 12\)} & M1 A1 & {\msNote{M1}{Divides 84 by the modulus of their determinant, or equates their scale factor times the area of \(T\) to 84.}} \\
& {\(\tfrac12\left|\det\begin{pmatrix}4 & k-1\\3 & -3\end{pmatrix}\right|\)\\[4pt] \(\phantom{A} = \tfrac12\left|-12 - 3(k-1)\right|\)} & M1 & {Finds the area of \(T\) in terms of \(k\) by any method, such as half a determinant of two side vectors or the shoelace formula. Condone sign slips.} \\
& {\(\phantom{A} = \tfrac32(k+3)\)} & A1 & {\(\frac32(k+3)\) or \(\frac32|k+3|\) oe.} \\
& {\(\tfrac32(k+3) = 12 \Rightarrow k = 5\)} & A1 & {\(k = 5\) only; A0 if \(k = -11\) is also given.} \\
\nopagebreak & & [6] & \\
\nopagebreak \SetCell[c=2]{r} Total for Question 3 & & 9 & \\
\end{msscheme}

\msQuestion{4}
\begin{msscheme}{
  hline{3-8,10-12,14} = {2-Z}{0.4pt, msGrid},
  hline{9,13,15} = {0.4pt, msGrid},
  row{15} = {font=\bfseries},
}
Question & Scheme & Marks & Notes \\
4(a) & {\(\begin{pmatrix}3\\1\end{pmatrix},\ \begin{pmatrix}2\\3\end{pmatrix}\)\\[4pt] \(\text{Area of } S = \left|\det\begin{pmatrix}3&2\\1&3\end{pmatrix}\right| = \ldots\)} & M1 & {Finds the area of \(S\) by any valid method, such as the determinant of two adjacent side vectors from one vertex. Condone a sign slip.} \\
& {\(\text{Area of } S = 7\)} & A1 & \\
& {\(\text{"}7\text{"}\,|\det\mathbf P| = 35\)} & M1 & {Equates their area of \(S\) times \(\det\mathbf P\), with or without the modulus, to 35, oe.} \\
& {\(\det\mathbf P = p^2 + 2p - 3 = \pm 5\)} & A1 & {Correct \(\det\mathbf P\) in any form, such as \((p+1)^2 - 4\), equated to 5 or to \(\pm 5\).} \\
& {\(p^2 + 2p - 8 = 0 \Rightarrow p = \ldots\)} & dM1 & {Solves their quadratic in \(p\) by factorising, the formula, completing the square or a calculator.} \\
& {\(p = 2\) \((p = -4 \text{ rejected})\)} & A1 & {\(p = 2\) only, with \(-4\) rejected. The case \(\det\mathbf P = -5\) need not be solved.} \\
\nopagebreak & & [6] & \\
4(b) & {\(\begin{pmatrix}0&1\\1&0\end{pmatrix}\)} & B1 & \\
& {\(\begin{pmatrix}2&0\\0&1\end{pmatrix}\)} & B1 & \\
& {\(\mathbf Q = \begin{pmatrix}2&0\\0&1\end{pmatrix}\begin{pmatrix}0&1\\1&0\end{pmatrix}\)\\[4pt] \(\phantom{\mathbf Q} = \begin{pmatrix}0&2\\1&0\end{pmatrix}\)} & B1ft & {ft their two matrices, stretch times reflection in that order, multiplied correctly.} \\
\nopagebreak & & [3] & \\
4(c) & {\(\mathbf R = \mathbf Q\mathbf P = \begin{pmatrix}0&2\\1&0\end{pmatrix}\begin{pmatrix}3&2\\2&3\end{pmatrix}\)\\[4pt] \(\phantom{\mathbf R} = \begin{pmatrix}4&6\\3&2\end{pmatrix}\)} & B1ft & {ft their \(\mathbf Q\) and their \(p\), with \(\mathbf Q\mathbf P\), in that order, multiplied correctly.} \\
\nopagebreak & & [1] & \\
\nopagebreak \SetCell[c=2]{r} Total for Question 4 & & 10 & \\
\end{msscheme}

\msQuestion{5}
\begin{msscheme}{
  hline{3,5,7-8,10-11,13,15-16} = {2-Z}{0.4pt, msGrid},
  hline{4,6,9,12,14,17} = {0.4pt, msGrid},
  row{17} = {font=\bfseries},
}
Question & Scheme & Marks & Notes \\
5(i) & {\(\mathbf{A}^2 = \begin{pmatrix} 1 & 0 \\ 0 & 1 \end{pmatrix}\)} & B1 & \\
\nopagebreak & & [1] & \\
5(ii) & {A reflection applied twice returns every point\\[4pt] to itself, so \(\mathbf{A}^2 = \mathbf{I}\) is consistent} & E1 & {States that reflecting twice leaves every point unchanged, or that a reflection is self-inverse, and links this to \(\mathbf{A}^2 = \mathbf{I}\).} \\
\nopagebreak & & [1] & \\
5(iii) & {\(-\tfrac12 x + \tfrac{\sqrt3}{2}y = x\)\\[4pt] \(\tfrac{\sqrt3}{2}x + \tfrac12 y = y\)} & M1 & {Sets \(\mathbf{A}\) times a general point \((x, y)\) equal to \((x, y)\), forming at least one equation. Condone slips in multiplying.} \\
& {\(y = \sqrt3\,x\)} & A1 & {oe, such as \(\sqrt3\,x - y = 0\).} \\
\nopagebreak & & [2] & \\
5(iv) & {Rotation about the origin} & B1 & {Needs the centre, the origin or \((0, 0)\).} \\
& {through \(120^\circ\) anticlockwise} & B1 & {Accept \(\frac{2\pi}{3}\), or \(240^\circ\) clockwise. The direction is required.} \\
\nopagebreak & & [2] & \\
5(v) & {\(\mathbf{BA} = \begin{pmatrix} -\frac12 & -\frac{\sqrt3}{2} \\ -\frac{\sqrt3}{2} & \frac12 \end{pmatrix}\)} & B1 & {The product in the order \(\mathbf{BA}\); \(\mathbf{AB}\) scores B0.} \\
\nopagebreak & & [1] & \\
5(vi) & {\(-\tfrac12 x - \tfrac{\sqrt3}{2}y = x\)\\[4pt] \(-\tfrac{\sqrt3}{2}x + \tfrac12 y = y\)} & M1 & {Sets their \(2 \times 2\) matrix from (v) times \((x, y)\) equal to \((x, y)\), forming at least one equation. Condone slips in multiplying.} \\
& {\(y = -\sqrt3\,x\)} & A1ft & {oe. ft their matrix from (v), where its invariant points form a line.} \\
\nopagebreak & & [2] & \\
\nopagebreak \SetCell[c=2]{r} Total for Question 5 & & 9 & \\
\end{msscheme}

\msQuestion{6}
\begin{msscheme}{
  hline{3-4,6-8,10-12} = {2-Z}{0.4pt, msGrid},
  hline{5,9,13} = {0.4pt, msGrid},
  row{13} = {font=\bfseries},
}
Question & Scheme & Marks & Notes \\
6(a) & {Stretch parallel to the \(y\)-axis, scale factor \(k\)} & B1 & {Names a stretch; its direction and scale factor are not needed, but B0 if a wrong one is given.} \\
& {then reflection in the line \(y = x\)} & B1 & {The reflection, after the stretch; its mirror line is not needed, but B0 if wrong. Reflection first scores B1 B1 only with the stretch parallel to the \(x\)-axis; otherwise B1 B0.} \\
\nopagebreak & & [2] & \\
6(b) & {\(\mathbf{M} = \begin{pmatrix} 0 & k \\ 1 & 0 \end{pmatrix}\), \(\det \mathbf{M} = -k\)\\[4pt] Area \(= k\)} & B1 & {Must be positive.} \\
& {\(\mathbf{M}^{-1} = -\frac{1}{k}\begin{pmatrix} 0 & -k \\ -1 & 0 \end{pmatrix}\)} & M1 & {Attempts the inverse of their \(\mathbf{M}\), by the determinant and adjugate, or as the inverses of the two matrices in reverse order. Condone one sign slip.} \\
& {\(\phantom{\mathbf{M}^{-1}} = \begin{pmatrix} 0 & 1 \\ \frac{1}{k} & 0 \end{pmatrix}\)} & A1 & \\
\nopagebreak & & [3] & \\
6(c) & {\(\begin{pmatrix} 0 & k \\ 1 & 0 \end{pmatrix}\begin{pmatrix} x \\ \frac{x}{\sqrt{k}} \end{pmatrix} = \ldots\)} & M1 & {Applies their \(\mathbf{M}\) to a general point of the line, such as \(\left(x, \frac{x}{\sqrt{k}}\right)\) or \((x, mx)\), or to its direction vector.} \\
& {\(= \begin{pmatrix} \sqrt{k}\,x \\ x \end{pmatrix}\)} & A1 & {oe, such as \(\left(\frac{kx}{\sqrt{k}}, x\right)\).} \\
& {\(X = \sqrt{k}\,x,\ Y = x \Rightarrow Y = \frac{X}{\sqrt{k}}\)\\[4pt] so the line is invariant} & A1~AG & {cso. Image shown to satisfy \(y = \frac{x}{\sqrt{k}}\), or to be a multiple of the direction vector, with a conclusion. Via \((x, mx)\), needs \(m^2 = \frac{1}{k}\) and \(m = \frac{1}{\sqrt{k}}\).} \\
\nopagebreak & & [3] & \\
\nopagebreak \SetCell[c=2]{r} Total for Question 6 & & 8 & \\
\end{msscheme}

\msQuestion{7}
\begin{msscheme}{
  hline{3} = {2-Z}{0.4pt, msGrid},
  hline{4} = {0.4pt, msGrid},
  row{4} = {font=\bfseries},
}
Question & Scheme & Marks & Notes \\
7 & {Compares with rotation about the \(y\)-axis:\\[4pt] \(\cos\theta = \frac{\sqrt3}{2}\), \(\sin\theta = \ldots\)} & M1 & {Uses the \(y\)-axis rotation matrix, with \(\sin\theta\) top right and \(-\sin\theta\) bottom left, to find the sign of \(\sin\theta\).} \\
& {\(\sin\theta = -\frac12\): Bilal ignored the sign\\[4pt] \(\theta = -30^\circ\) (or \(330^\circ\))} & A1 & {Error named as the sign of \(\sin\theta\), and \(\theta = -30^\circ\) or \(330^\circ\) oe, or \(30^\circ\) clockwise viewed from the positive \(y\)-axis.} \\
\nopagebreak \SetCell[c=2]{r} Total for Question 7 & & 2 & \\
\end{msscheme}

\msQuestion{8}
\begin{msscheme}{
  hline{3,5-6,8,10} = {2-Z}{0.4pt, msGrid},
  hline{4,7,9,11} = {0.4pt, msGrid},
  row{11} = {font=\bfseries},
}
Question & Scheme & Marks & Notes \\
8(a) & {\(\mathbf{A}^5 = \begin{pmatrix}0 & 0 & 1\\1 & 0 & 0\\0 & 1 & 0\end{pmatrix}\)} & B1 & \\
\nopagebreak & & [1] & \\
8(b) & {Rotation about the line \(x = y = z\)} & B1 & {Rotation, with axis \(x = y = z\) oe, such as the direction \(\mathbf{i} + \mathbf{j} + \mathbf{k}\).} \\
& {through \(120^\circ\), clockwise viewed from \((1, 1, 1)\) towards the origin} & B1 & {Angle with a correct sense, oe such as anticlockwise viewed from the origin towards \((1, 1, 1)\), or \(240^\circ\) anticlockwise from \((1, 1, 1)\). B0 without a sense.} \\
\nopagebreak & & [2] & \\
8(c) & {\(\mathbf{B} = \begin{pmatrix}0 & 1 & 0\\1 & 0 & 0\\0 & 0 & 1\end{pmatrix}\)} & B1 & \\
\nopagebreak & & [1] & \\
8(d) & {\(Q'\) is \((-1, 4, 2)\)} & B1ft & {ft their \(\mathbf{B}\) applied to \((4, -1, 2)\). Condone a column vector.} \\
\nopagebreak & & [1] & \\
\nopagebreak \SetCell[c=2]{r} Total for Question 8 & & 5 & \\
\end{msscheme}

\msQuestion{9}
\begin{msscheme}{
  hline{3,5-6,8,10-11,13-15,17-18} = {2-Z}{0.4pt, msGrid},
  hline{4,7,9,12,16,19} = {0.4pt, msGrid},
  row{19} = {font=\bfseries},
}
Question & Scheme & Marks & Notes \\
9(a) & {\(\alpha = 225\)} & B1 & {Or any \(225 + 360n\), such as \(-135\).} \\
\nopagebreak & & [1] & \\
9(b) & {\(225k = 360n \Rightarrow 5k = 8n \Rightarrow k = \ldots\)} & M1 & {Sets \(k\) times their \(\alpha\) equal to a multiple of 360, or to 360 times an integer, and attempts the least \(k\).} \\
& {\(k = 8\)} & A1 & \\
\nopagebreak & & [2] & \\
9(c) & {\(\mathbf{N} = \begin{pmatrix}-1&0&0\\0&1&0\\0&0&1\end{pmatrix}\)} & B1 & \\
\nopagebreak & & [1] & \\
9(d) & {\(\mathbf{N}\mathbf{M}\begin{pmatrix}1\\2\\3\end{pmatrix} = \mathbf{N}\begin{pmatrix}-2\sqrt2\\2\\-\sqrt2\end{pmatrix} = \ldots\)} & M1 & {Multiplies \((1, 2, 3)\) by \(\mathbf{M}\), then by their reflection matrix \(\mathbf{N}\), or by their \(\mathbf{NM}\), in that order. Condone slips.} \\
& {\(\phantom{\mathbf{N}\mathbf{M}} = \begin{pmatrix}2\sqrt2\\2\\-\sqrt2\end{pmatrix} \Rightarrow B = (2\sqrt2, 2, -\sqrt2)\)} & A1~AG & {Correct products with \(\mathbf{M}\) and the correct \(\mathbf{N}\), surds simplified, leading without error to the printed coordinates.} \\
\nopagebreak & & [2] & \\
9(e) & {\(|\overrightarrow{OA}| = \sqrt{14}\), \(|\overrightarrow{OB}| = \sqrt{8 + 4 + 2} = \sqrt{14}\)} & B1 & \\
& {\(\cos\angle AOB = \frac{2\sqrt2 + 4 - 3\sqrt2}{\sqrt{14}\sqrt{14}}\)} & M1 & {Divides a three-term scalar product of \(\overrightarrow{OA}\) and \(\overrightarrow{OB}\) by their magnitudes, or uses the cosine rule with their \(|AB|\). Condone slips.} \\
& {\(\phantom{\cos\angle AOB} = \frac{4 - \sqrt2}{14} \Rightarrow \angle AOB = 79.4^\circ\)} & A1~AG & {Correct \(\cos\angle AOB = \frac{4-\sqrt2}{14}\) oe, leading without error to \(79.4^\circ\).} \\
\nopagebreak & & [3] & \\
9(f) & {\(\text{Area} = \frac12\sqrt{14}\sqrt{14}\sin 79.4^\circ\)} & M1 & {Uses \(\frac12|\overrightarrow{OA}||\overrightarrow{OB}|\sin\angle AOB\) with their magnitudes and \(79.4^\circ\), a more accurate angle or the exact sine.} \\
& {\(\phantom{\text{Area}} = 6.88\)} & A1 & \\
\nopagebreak & & [2] & \\
\nopagebreak \SetCell[c=2]{r} Total for Question 9 & & 11 & \\
\end{msscheme}

\msQuestion{10}
\begin{msscheme}{
  hline{3-5,7} = {2-Z}{0.4pt, msGrid},
  hline{6,8} = {0.4pt, msGrid},
  row{8} = {font=\bfseries},
}
Question & Scheme & Marks & Notes \\
10(a) & {\(\mathbf{R} = \begin{pmatrix}1&0&0\\0&\cos\theta&-\sin\theta\\0&\sin\theta&\cos\theta\end{pmatrix}\)\\[4pt] \(\Rightarrow \mathbf{R}^{-1} = \begin{pmatrix}1&0&0\\0&\cos\theta&\sin\theta\\0&-\sin\theta&\cos\theta\end{pmatrix}\)} & B1 & {Correct \(\mathbf{R}\) or \(\mathbf{R}^{-1}\), stated or used.} \\
& {\(\mathbf{A} = \mathbf{R}^{-1}\mathbf{B} = \ldots\)} & M1 & {Uses \(\mathbf{A} = \mathbf{R}^{-1}\mathbf{B}\), not \(\mathbf{B}\mathbf{R}^{-1}\), with their \(\mathbf{R}^{-1}\) of a rotation matrix's form, and attempts the product. Condone slips in multiplying.} \\
& {\(\mathbf{A} = \begin{pmatrix}1&0&0\\0&0&\cos^2\theta+\sin^2\theta\\0&\sin^2\theta+\cos^2\theta&0\end{pmatrix}\)\\[4pt] \(\phantom{\mathbf{A}} = \begin{pmatrix}1&0&0\\0&0&1\\0&1&0\end{pmatrix}\), so no \(\theta\)} & A1~AG & {Correct product simplified by \(\sin^2\theta+\cos^2\theta=1\) to a matrix free of \(\theta\), with a conclusion that \(\mathbf{A}\) is independent of \(\theta\).} \\
\nopagebreak & & [3] & \\
10(b) & {Reflection in the plane \(y = z\)} & B1 & {Needs reflection and the plane, \(y = z\) oe.} \\
\nopagebreak & & [1] & \\
\nopagebreak \SetCell[c=2]{r} Total for Question 10 & & 4 & \\
\end{msscheme}

\msQuestion{11}
\begin{msscheme}{
  hline{3-4,6,8-9,11-16,18} = {2-Z}{0.4pt, msGrid},
  hline{5,7,10,17,19} = {0.4pt, msGrid},
  row{19} = {font=\bfseries},
}
Question & Scheme & Marks & Notes \\
11(a) & {Rotation about the origin} & B1 & \\
& {through angle \(\theta\) anticlockwise} & B1 & {Accept clockwise through \(2\pi - \theta\).} \\
\nopagebreak & & [2] & \\
11(b)(i) & {\(\det(\mathbf{M}) = \cos^2\theta + \sin^2\theta = 1\)} & B1 & \\
\nopagebreak & & [1] & \\
11(b)(ii) & {\(|\det(\mathbf{M})| = 1 \Rightarrow\) area is preserved} & E1 & {States that area is preserved, or the area scale factor is 1, linked to the determinant.} \\
& {\(\det(\mathbf{M}) > 0 \Rightarrow\) orientation is preserved} & E1 & {States that orientation is preserved, or that \(T\) involves no reflection, linked to the determinant being positive.} \\
\nopagebreak & & [2] & \\
11(c) & {\(n = 1\): \(\cos 1\theta = \cos\theta\), \(\sin 1\theta = \sin\theta\)\\[4pt] \(\Rightarrow \mathbf{M}^1 = \mathbf{M}\), so true for \(n = 1\)} & B1 & {Formula evaluated at \(n = 1\) and seen to give \(\mathbf{M}\).} \\
& {Assumes true for \(n = k\):\\[4pt] \(\mathbf{M}^{k+1} = \mathbf{M}^k\mathbf{M}\)} & M1 & {Assumes the result for \(n = k\) and writes \(\mathbf{M}^{k+1}\) as \(\mathbf{M}^k\mathbf{M}\) or \(\mathbf{M}\mathbf{M}^k\). Not if the result for \(n = k\) is stated as fact.} \\
& {\(\begin{pmatrix}\cos k\theta & -\sin k\theta \\ \sin k\theta & \cos k\theta\end{pmatrix}\begin{pmatrix}\cos\theta & -\sin\theta \\ \sin\theta & \cos\theta\end{pmatrix} = \ldots\)} & M1 & {Multiplies the assumed form of \(\mathbf{M}^k\) and \(\mathbf{M}\), in either order, each entry a sum of two products. Condone sign slips.} \\
& {\(\begin{pmatrix} c & -s \\ s & c\end{pmatrix}\), where\\[4pt] \(c = \cos k\theta\cos\theta - \sin k\theta\sin\theta\)\\[4pt] \(s = \sin k\theta\cos\theta + \cos k\theta\sin\theta\)} & A1 & {All four entries correct, unsimplified.} \\
& {\(\Rightarrow \mathbf{M}^{k+1} = \begin{pmatrix}\cos(k+1)\theta & -\sin(k+1)\theta \\ \sin(k+1)\theta & \cos(k+1)\theta\end{pmatrix}\)} & A1 & {cso. All four entries in the form for \(n = k + 1\), using the addition formulae.} \\
& {True for \(n = 1\); true for \(n = k\)\\[4pt] \(\Rightarrow\) true for \(n = k + 1\)\\[4pt] \(\Rightarrow\) true for all positive integers \(n\)} & E1 & {Concludes with all three points, true for \(n = 1\), true for \(n = k\) implies true for \(n = k + 1\), and so true for all positive integers \(n\).} \\
\nopagebreak & & [6] & \\
11(d) & {Rotation through \(n\theta\) anticlockwise about the origin} & B1ft & {ft their description in (a) with \(\theta\) replaced by \(n\theta\). Needs the centre, angle and sense. Condone the angle not reduced modulo \(2\pi\).} \\
\nopagebreak & & [1] & \\
\nopagebreak \SetCell[c=2]{r} Total for Question 11 & & 12 & \\
\end{msscheme}

\end{document}
